%----------------------------------------------------------------------------------------
%	PACKAGES AND OTHER DOCUMENT CONFIGURATIONS
%----------------------------------------------------------------------------------------

\documentclass[
12pt, % The default document font size, options: 10pt, 11pt, 12pt
oneside, % Two side (alternating margins) for binding by default, uncomment to switch to one side
chapterinoneline,% Have the chapter title next to the number in one single line
english, % ngerman for German
onehalfspacing, % Single line spacing, alternatives: onehalfspacing or doublespacing
%draft, % Uncomment to enable draft mode (no pictures, no links, overfull hboxes indicated)
nolistspacing, % If the document is onehalfspacing or doublespacing, uncomment this to set spacing in lists to single
%liststotoc, % Uncomment to add the list of figures/tables/etc to the table of contents
%toctotoc, % Uncomment to add the main table of contents to the table of contents
%parskip, % Uncomment to add space between paragraphs
%nohyperref, % Uncomment to not load the hyperref package
headsepline, % Uncomment to get a line under the header
]{MastersDoctoralThesis} % The class file specifying the document structure

\usepackage{cmap}

\usepackage[utf8]{inputenc} % Required for inputting international characters
\usepackage[T2A]{fontenc} % Output font encoding for international characters

\usepackage[backend=bibtex,style=numeric,natbib=true]{biblatex} % Use the bibtex backend with the authoryear citation style (which resembles APA)

\addbibresource{references.bib} % The filename of the bibliography

\usepackage[autostyle=true]{csquotes} % Required to generate language-dependent quotes in the bibliography

\usepackage{wrapfig}
\usepackage{float}

\usepackage{caption}
\captionsetup[table]{name=Таблица}
\captionsetup[figure]{name=Фигура}

\usepackage{longtable}

\usepackage{indentfirst}


\usepackage{pgfplots}
\pgfplotsset{
	width=12.5cm,
    compat=1.9,
    height=7.5cm,
	xlabel={Номер на итерация},
    ylabel={Точност в \%},
    xmin=1, xmax=25,
    ymin=0, ymax=100,
    xtick={1,5,10,15,20,25},
    ytick={0,20,40,60,80,100},
    legend style={at={(0.5,-0.2)},anchor=north},
    legend cell align=left,
    ymajorgrids=true,
    grid style=dashed,}
    
\usepackage{microtype}

%----------------------------------------------------------------------------------------
%	MARGIN SETTINGS
%----------------------------------------------------------------------------------------

\geometry{
	paper=a4paper, % Change to letterpaper for US letter
	inner=2.5cm, % Inner margin
	outer=3.8cm, % Outer margin
	bindingoffset=2cm, % Binding offset
	top=1.5cm, % Top margin
	bottom=1.5cm, % Bottom margin
    head=29pt,
	%showframe,% show how the type block is set on the page
}

%----------------------------------------------------------------------------------------
%	THESIS INFORMATION
%----------------------------------------------------------------------------------------

\hypersetup{pdftitle=Автоматично разделяне на големи количества изображения на класове с използване на дълбоко учене (deep learning)} % Set the PDF's title to your title
\hypersetup{pdfauthor=Валентин Змийчаров} % Set the PDF's author to your name

\begin{document}
\sloppy

%\frontmatter % Use roman page numbering style (i, ii, iii, iv...) for the pre-content pages

\mainmatter % Begin numeric (1,2,3...) page numbering

\pagestyle{thesis} % Default to the plain heading style until the thesis style is called for the body content

%----------------------------------------------------------------------------------------
%	TITLE PAGE
%----------------------------------------------------------------------------------------

\begin{titlepage}
\begin{center}

\Large \textbf{Софийски университет \\ „Св. Климент Охридски”}\\[0cm]
\Large Дипломна работа\\[0cm]
\Large \textit{Катедра „Софтуерни технологии”}\\[0.2cm]

\includegraphics{Figures/logo_su.jpg}
\includegraphics{Figures/logo_katedra.jpg}\\[0.1cm]

\HRule \\[0.1cm] % Horizontal line
{\huge \bfseries Прилагане на методи от машинния превод за изграждане на диалогова система за отговаряне на често задавани въпроси
	\par}\vspace{0.1cm} % Thesis title
\HRule \\[0.2cm] % Horizontal line
 
\begin{flushleft} \large
Дипломант:\\ \textbf{Мартин Боянов}\\
Специалност:\\ \textbf{Извличане на информация и откриване на знания}\\
Факултетен номер:\\ \textbf{25187} \\
Научни ръководители:\\ \textbf{Д-р Преслав Наков, Проф. Иван Койчев}\\[0.2cm]
\end{flushleft}

{\large София 2017г.}\\
 
\vfill
\end{center}
\end{titlepage}

%----------------------------------------------------------------------------------------
%	LIST OF CONTENTS/FIGURES/TABLES PAGES
%----------------------------------------------------------------------------------------

\renewcommand{\contentsname}{Съдържание}
\tableofcontents % Prints the main table of contents

%----------------------------------------------------------------------------------------
%	THESIS CONTENT - CHAPTERS
%----------------------------------------------------------------------------------------


% \pagestyle{thesis} % Return the page headers back to the "thesis" style

\chapter{Увод}

\label{Chapter1}

%----------------------------------------------------------------------------------------

\section{Актуалност на проблема и мотивация}
Бързо разрастващата се сфера на мигновените съобщения (instant-messaging) поражда голям интерес за създаване на интелигентни системи за водене на диалог, още известни като чатботове.\\
Изграждането на такова приложение е трудна задача, поради липсата на структура в текста. Трудността се засилва и от факта, че при чат приложенията правописът и граматиката не са приоритет за потребителите. При зададен въпрос от потребител, системата трябва да генерира смислен и релевантен отговор. \\
Подобна задача се решава при машинния превод. При дадена поредица от думи, трябва да се генерира друга поредица от думи. Целта на тази дипломна работа е да се адаптира модел от машинния превод за изграждане на чатбот.



\section{Цел и задачи на дипломната работа}
Основната цел на тази работа е изследване на възможностите за изграждане на диалогова система базирана на модела seq2seq. Този модел е в основата на съвременните системи за машинен превод, използващи изкуствени невронни мрежи. Ще се изследват и възможностите за интегриране на модела в система за автоматично отговаряне на въпроси.

Задачи, произтичащи от целта:
\begin{itemize}
	\item Обзор на подходите за решаване на този вид задачи
	\item Набавяне на набор от данни - Ubuntu Dataset Corpus и QatarLiving
	\item Предварителна обработка на данните
	\item Имплементация и интеграция на seq2seq модела
	\item Анализ на процеса на машинно самообучение
	\item Тестване, настройка и оценяване на работата на системата
	\item Анализ на резултатите
	\item Интеграция в системата за автоматично отговаряне на въпроси
	\item Разработване на интерактивно демонстративно приложение
\end{itemize}


\section{Структура на дипломната работа}



\chapter{Преглед на предметната област} % Main chapter title

\label{Chapter2} 

%----------------------------------------------------------------------------------------

\section{Въведение}
Изграждането на отворена диалогова система е предизвикателна задача, с която са се сблъсквали много екипи. През годините са опитвани различни подходи за изграждане на подобна система и в тази глава ще се запознаем с основните методи и тяхната еволюция. 
Благодарение на прогреса в сферата на машинното самообучение и на нарастващата изчислителна мощ бяха постигнати подобрения в задачата за автоматичен машинен превод. Ще обобщим тяхната архитектура и ще коментираме приложимостта на подобен модел за изграждане на диалогова система.
Ще разгледаме и проблема за отговаряне на често задавани въпроси и връзката му с диалоговите системи


https://research.fb.com/the-long-game-towards-understanding-dialog/

\section{Диалогови системи: мотивация и история}

-мотивация
-ELIZA
-PARRY?
-ALICE
-games
-assistants

\section{Машинен превод}
- statistical machine translation
- neural machine translation
- seq2seq


\section{Отговаряне на често задавани въпроси}
 
\chapter{Използвани технологии и платформи}

\label{Chapter3}

%----------------------------------------------------------------------------------------

\section{intro}

\section{tensorflow}
\section{python}
\section{kelp}
\section{BPE}
\section{that person with the chatbot design}
\section{Data}
\subsection{UDC}
\subsection{QL}


\chapter{Реализация, експерименти и анализ на резултатите}

\label{Chapter4}

%----------------------------------------------------------------------------------------

\section{Реализация}







 
\chapter{Заключение}

\label{Chapter5}

%----------------------------------------------------------------------------------------

\section{Обобщение на изпълнението на началните цели}


 

%----------------------------------------------------------------------------------------
%	BIBLIOGRAPHY
%----------------------------------------------------------------------------------------

\printbibliography[title={Референции},heading=bibnumbered]

%----------------------------------------------------------------------------------------

%----------------------------------------------------------------------------------------
%	APPLICATIONS
%----------------------------------------------------------------------------------------

\chapter{Приложения}

\label{Applications}

%----------------------------------------------------------------------------------------










%---------------------------------------------------------------------------------------- 

\end{document}  
